\chapter{群论}\label{ch:b3:groups}

本科二年级分册把群当作记账工具：拉格朗日
定理、循环群、对称群及其符号。本章把群论
变为一种\emph{方法}。其核心是群
\emph{作用}于集合这一概念：计数轨道与不动点，
得出类方程、柯西定理以及三个西罗
定理——它们是有限群论中基本的
局部到整体原理。随后我们学习组装群（直积与
半直积），也学会拆解群（合成列、
可解群），并证明一个定理，它将在
\cref{ch:b3:galois}中终结一个存在了三百年的关于
多项式方程的问题：交错群 $A_n$ 当 $n
\geq 5$ 时是单群。

\section{商群与同构定理}

本节始终设 $G$ 为乘法记号写出的群，$e$
回顾本科二年级分册中的内容：子群、陪集 $gH$、
拉格朗日定理（$\abs G = [G:H]\,\abs H$，对有限群 $G$），
元素的阶、循环群，以及对称群 $S_n$
及其符号同态 $\varepsilon \colon S_n \to \{\pm1\}$。

\begin{definition}\label{def:b3:groups:normal}
称 $N$ 为 $G$ 的\emph{正规子群}\index{正规子群}
（记作 $N \trianglelefteq G$），当 $gNg^{-1} = N$ 对每个 $g
\in G$ 成立时——等价地，左陪集与右陪集重合：$gN =
Ng$ 对所有 $g$ 成立。
\end{definition}

\begin{theorem}[商群]\label{thm:b3:groups:quotient}
设 $N \trianglelefteq G$。陪集集合 $G/N$ 在
乘法 $(gN)(hN) = ghN$ 下构成群，
\emph{商群}\index{商群}，而
\emph{典范投影} $\pi \colon G \to G/N$，$g \mapsto gN$，
是满同态，其核为 $N$。反之，群同态的每个核
都是正规的：正规子群恰好就是
核。
\end{theorem}

\begin{proof}
良定性正是关键所在。若 $gN = g'N$，$hN = h'N$，
写 $g' = gn$，$h' = hm$，其中 $n, m \in N$。则 $g'h' = gnhm =
gh\,(h^{-1}nh)\,m \in ghN$，由正规性知 $h^{-1}nh \in N$，
故陪集之积不依赖于代表元的选取。
结合律、单位元 $eN = N$ 以及逆元 $(gN)^{-1} = g^{-1}N$
都继承自 $G$。显然 $\pi$ 是满同态，且
$\pi(g) = N \iff g \in N$。

若 $f \colon G \to H$ 为同态，$k \in \ker f$，则
$f(gkg^{-1}) = f(g)f(k)f(g)^{-1} = e$：核是正规的。
\end{proof}

\begin{theorem}[泛性质；第一同构定理]
\label{thm:b3:groups:firstiso}
设 $f \colon G \to H$ 为同态，$N \trianglelefteq G$ 且
$N \subseteq \ker f$。则存在唯一的同态 $\bar f \colon G/N
\to H$ 使得 $f = \bar f \circ \pi$。特别地，取 $N = \ker
f$：\index{同构定理}
\[
G/\ker f \;\xrightarrow{\;\sim\;}\; \operatorname{im} f,
\qquad gN \mapsto f(g).
\]
\end{theorem}

\begin{proof}
唯一性：$\bar f(gN)$ 必须等于 $f(g)$。存在性：若 $gN = g'N$，
则 $g^{-1}g' \in N \subseteq \ker f$，故 $f(g) = f(g')$，于是 $\bar
f(gN) = f(g)$ 是良定的；由于 $f$ 是同态，此映射亦为同态。对
$N = \ker f$：$\bar f$ 是单射，因为 $\bar f(gN) = e$ 表示
$g \in \ker f$，即 $gN = N$；其像与 $f$ 的像相同。
\end{proof}

\begin{theorem}[第二、第三同构定理]
\label{thm:b3:groups:secondiso}
设 $H \leq G$，$N \trianglelefteq G$。
\begin{enumerate}
  \item $HN = \{hn : h \in H,\, n \in N\}$ 是子群，$N
        \trianglelefteq HN$，$H \cap N \trianglelefteq H$，且
        \[
        H/(H \cap N) \;\cong\; HN/N .
        \]
  \item 若又有 $N \subseteq K \trianglelefteq G$，则 $K/N
        \trianglelefteq G/N$，且 $(G/N)\big/(K/N) \cong G/K$。
\end{enumerate}
\end{theorem}

\begin{proof}
(1) $HN$ 是子群：$(hn)(h'n') = hh'\,(h'^{-1}nh')n' \in HN$，
且 $(hn)^{-1} = h^{-1}(hn^{-1}h^{-1}) \in HN$，利用
$N$ 的正规性。复合 $H \hookrightarrow HN \xrightarrow{\pi} HN/N$：该
同态是满的（$hnN = hN$），其核为 $\{h \in H : h \in
N\} = H \cap N$；应用\cref{thm:b3:groups:firstiso}。

(2) 投影 $G/N \to G/K$，$gN \mapsto gK$，
由 $N \subseteq K$ 良定，满射，核为 $K/N$；再次应用
\cref{thm:b3:groups:firstiso}。
\end{proof}

\begin{theorem}[对应定理]\label{thm:b3:groups:correspondence}
设 $N \trianglelefteq G$。映射 $H \mapsto H/N$ 是
$G$ 中包含 $N$ 的子群与 $G/N$ 的子群之间的双射，
保持包含关系、指数与正规性（双向均成立）。
\end{theorem}

\begin{proof}
其逆映射为 $\bar H \mapsto \pi^{-1}(\bar H)$。两个映射都把
子群送到子群，且互为逆：$\pi^{-1}(H/N) =
HN = H$（因 $N \subseteq H$），且 $\pi(\pi^{-1}(\bar H)) = \bar H$，
由 $\pi$ 的满射性。包含关系显然保持；
$[G:H] = [G/N : H/N]$，因为 $gH \mapsto (gN)(H/N)$ 是
陪集空间之间良定的双射；又 $gHg^{-1} = H$ 对所有 $g$ 成立，
当且仅当 $(gN)(H/N)(gN)^{-1} = H/N$ 对所有 $gN$ 成立，
这同样由 $\pi$ 的满射性得到。
\end{proof}

\begin{example}\label{ex:b3:groups:firstiso}
$\varepsilon \colon S_n \to \{\pm 1\}$ 给出 $S_n/A_n \cong
\{\pm1\}$；$\det \colon GL_n(K) \to K^\times$ 给出
$GL_n(K)/SL_n(K) \cong K^\times$；$t \mapsto \eu^{2\iu\pi t}$ 给出
$\R/\Z \cong \mathbb U$，即圆群。第一同构定理是
在实践中\emph{计算}商的方法：
找一个具有正确核的满射。
\end{example}

\begin{method}\label{met:b3:groups:normal}
为证明 $N \trianglelefteq G$，按优雅程度由高到低：
将 $N$ 表示为定义在 $G$ 上的某个态射的核；
验证 $gNg^{-1} \subseteq N$ 对所有 $g$ 成立（这已足够：对
$g^{-1}$ 应用此式并共轭，即得反向包含）；验证
$N$ 是若干共轭类之并；或注意到 $[G:N] = 2$
（于是 $gN = Ng$ 必然成立 --- \cref{exo:b3:groups:1}）。
\end{method}

\section{群作用}

\begin{definition}\label{def:b3:groups:action}
称\emph{群作用}\index{群作用}为 $G$ 在集合 $X$ 上的一个
态射 $\varphi \colon G \to \mathfrak{S}(X)$，映入
$X$ 的双射群；记 $g \cdot x$ 为 $\varphi(g)(x)$。
等价地，即一个映射 $G \times X \to X$，满足 $e \cdot x = x$ 且 $g
\cdot (h \cdot x) = (gh) \cdot x$。$x$ 的\emph{轨道}\index{轨道}
为 $\mathcal O_x = \{g \cdot x : g \in G\}$，其
\emph{稳定子}\index{稳定子}为子群 $G_x = \{g :
g\cdot x = x\}$，而 $X^G = \{x : \forall g,\ g \cdot x = x\}$ 是
\emph{不动点}的集合。称作用为\emph{传递的}，若
恰有一个轨道，称为\emph{忠实的}，若 $\varphi$ 为
单射；称为\emph{自由的}，若所有稳定子都是平凡的。
\end{definition}

\begin{example}\label{ex:b3:groups:actions}
以下五个作用贯穿有限群论的始终：
\begin{enumerate}
  \item $G$ 通过\emph{左平移} $g \cdot x = gx$ 作用于自身：
        自由且传递。
  \item $G$ 通过\emph{共轭}作用于自身，$g \cdot x = gxg^{-1}$：
        轨道即\emph{共轭类}\index{共轭类}，
        稳定子为\emph{中心化子}
        \index{中心化子} $Z_G(x) = \{g : gx = xg\}$，不动点为
        \emph{中心}\index{中心} $Z(G)$。
  \item $G$ 在陪集空间 $G/H$ 上，通过 $g \cdot xH = gxH$ 作用：
        传递的，其中陪集 $H$ 的稳定子等于 $H$。
        每个传递作用都具有这种形式
        （见\cref{exo:b3:groups:8}）。
  \item $G$ 通过共轭作用于其子群全体：
        $H$ 的稳定子为\emph{正规化子}\index{正规化子} $N_G(H) =
        \{g : gHg^{-1} = H\}$，它是 $G$ 中使 $H$
        为正规子群的最大子群。
  \item $S_n$ 作用于 $\intint{1}{n}$：一切例子之母。
\end{enumerate}
\end{example}

\begin{theorem}[轨道--稳定子]\label{thm:b3:groups:orbitstab}
映射 $gG_x \mapsto g \cdot x$ 是良定的双射
$G/G_x \to \mathcal O_x$。特别地，当 $G$
有限时，\index{轨道--稳定子定理}
\[
\abs{\mathcal O_x} = [G : G_x] \quad\text{整除 } \abs G,
\]
并且，轨道划分 $X$（它们是
等价关系 $x \sim y \iff y \in \mathcal O_x$ 的等价类），
\[
\abs X = \sum_{i} \,[G : G_{x_i}]
\qquad (x_i\colon \text{每个轨道取一点}).
\]
\end{theorem}

\begin{proof}
良定且单射：$gG_x = hG_x \iff h^{-1}g \in G_x \iff
h^{-1}g \cdot x = x \iff g \cdot x = h \cdot x$；此链须从两个方向读取。
满射性即轨道的定义。
计数的结论由拉格朗日定理及
$X$ 按轨道的分拆得出。
\end{proof}

\begin{corollary}[类方程]\label{cor:b3:groups:classeq}
设 $G$ 为有限群，在每个
元素多于一个的共轭类中各取一个代表元 $x_i$：\index{类方程}
\[
\abs G = \abs{Z(G)} + \sum_i \,[G : Z_G(x_i)],
\qquad\text{每个 } [G : Z_G(x_i)] > 1 \text{ 且整除 } \abs G.
\]
\end{corollary}

\begin{proof}
对共轭作用应用\cref{thm:b3:groups:orbitstab}：
单点轨道恰为 $Z(G)$ 中的元素。
\end{proof}

\begin{theorem}[$p$-群的不动点]\label{thm:b3:groups:pfixed}
设 $p$ 为质数。称\emph{$p$-群}\index{$p$-群}为
阶为 $p$ 的幂的有限群。若 $p$-群 $G$ 作用于
有限集 $X$ 上，则
\[
\abs{X^G} \equiv \abs X \pmod p .
\]
由此可得：非平凡的 $p$-群的中心非平凡，且
每个阶为 $p^2$ 的群都是阿贝尔群。
\end{theorem}

\begin{proof}
每个轨道的基数为 $[G:G_x]$，是 $p$ 的幂；此幂
恰在不动点处为 $1$，其余情形可被 $p$ 整除。对各轨道求和
得同余式。关于中心，$G$ 对自身的共轭作用
的不动点集满足 $X^G = Z(G)$，故 $\abs{Z(G)} \equiv
\abs G \equiv 0 \pmod p$，且 $Z(G) \ni e$ 迫使 $\abs{Z(G)} \geq
p$。阶为 $p^2$ 的群：若 $Z(G) \neq G$，则 $\abs{Z(G)} = p$，且
$G/Z(G)$ 为 $p$ 阶循环群，这迫使 $G$ 为阿贝尔群
（见\cref{exo:b3:groups:2}） --- 矛盾。
\end{proof}

\begin{theorem}[柯西]\label{thm:b3:groups:cauchy}
若素数 $p$ 整除 $\abs G$，则 $G$ 含有一个
阶为 $p$ 的元素。\index{柯西定理}
\end{theorem}

\begin{proof}[证明（麦凯）]
设 $X = \{(g_1, \dots, g_p) \in G^p : g_1 g_2 \cdots g_p = e\}$。
任取 $g_1, \dots, g_{p-1}$ 即唯一确定 $g_p$：$\abs X =
\abs G^{p-1}$，可被 $p$ 整除。循环群 $\Z/p\Z$ 作用于
$X$ 上，即循环移位 $(g_1, \dots, g_p) \mapsto (g_2, \dots, g_p,
g_1)$ --- 这保持 $X$ 不变，因为 $g_2 \cdots g_p g_1 =
g_1^{-1}(g_1 \cdots g_p)g_1 = e$。由\cref{thm:b3:groups:pfixed}
可知 $\abs{X^{\Z/p\Z}} \equiv \abs X \equiv 0 \pmod p$。不动点恰为
常值元组 $(g, \dots, g)$，其中 $g^p = e$；元组 $(e,
\dots, e)$ 为其中之一，因此其个数至少为 $p$，从而
存在 $g \neq e$ 满足 $g^p = e$，其阶恰为 $p$。
\end{proof}

\begin{theorem}[凯莱]\label{thm:b3:groups:cayley}
每个 $n$ 阶群均可嵌入 $S_n$。\index{凯莱定理}
\end{theorem}

\begin{proof}
左平移 $\varphi \colon G \to \mathfrak S(G) \cong S_n$ 是
同态；$\varphi(g) = \mathrm{id}$ 推出 $g = ge = e$，
故它是忠实的。
\end{proof}

\begin{method}\label{met:b3:groups:counting}
不动点计数是有限群论的通用首选
手段。要证明某物\emph{存在}（例如中心元，
阶为 $p$ 的元素、正规子群、不动点），可取一个恰当的群
作用于一个恰当的有限集上，然后比较
$\abs{X^G}$ 与 $\abs X$ 模 $p$，或令轨道大小整除
群的阶。
\cref{thm:b3:groups:pfixed,thm:b3:groups:cauchy} 及下文三个西罗定理的证明
都是同一思想的五种变体。
\end{method}

\section{西罗定理}

拉格朗日定理指出，子群的阶整除 $\abs G$；
其逆不成立（$A_4$ 阶为 $12$，不含阶为
$6$ 的子群 --- 见\cref{exo:b3:groups:1}）。西罗定理挽回了
素数幂情形下的逆命题，其计数部分是
构造正规子群最精细的一般工具。

\begin{definition}\label{def:b3:groups:sylow}
记 $\abs G = p^a m$，其中 $p \nmid m$。称 \emph{西罗
$p$-子群}\index{西罗子群}为 $G$ 的阶为
$p^a$ 的子群 --- 阶可能最大的 $p$-子群。
西罗 $p$-子群在 $G$ 中的个数记为 $n_p$。
\end{definition}

\begin{lemma}\label{lem:b3:groups:binom}
若 $\abs G = p^a m$，其中 $p \nmid m$，则 $\dbinom{p^a m}{p^a}
\equiv m \pmod p$。
\end{lemma}

\begin{proof}
在 $\mathbb F_p[X]$ 中，由“新生梦” $(1+X)^p = 1 + X^p$（
系数 $\binom pk$，$0<k<p$，皆可被 $p$ 整除：$p$
整除 $\frac{p!}{k!(p-k)!}$ 的分子而不整除其
分母）迭代可得 $(1+X)^{p^a} = 1 + X^{p^a}$，于是
\[
(1+X)^{p^a m} = \bigl(1 + X^{p^a}\bigr)^m
= \sum_{k=0}^{m} \binom mk X^{k p^a} \quad\text{在 } \mathbb
F_p[X].
\]
比较 $X^{p^a}$ 的系数：左边为 $\binom{p^a
m}{p^a} \bmod p$，右边为 $\binom m1 = m$。
\end{proof}

\begin{theorem}[西罗 I：存在性]\label{thm:b3:groups:sylow1}
对每个素数 $p$，西罗 $p$-子群在 $G$ 中均存在。
\end{theorem}

\begin{proof}[证明（维兰特）]
设 $\Omega$ 为 \emph{子集}全体，这些子集取自 $G$，
且基数为 $p^a$；$G$ 通过左平移作用于 $\Omega$，$g \cdot S = gS$。
由\cref{lem:b3:groups:binom}，$\abs\Omega = \binom{p^a m}{p^a}
\equiv m \not\equiv 0 \pmod p$，故存在某轨道 $\mathcal O_S$，
其大小与 $p$ 互素（若 $p$ 整除每个轨道的大小，则它也整除
$\abs\Omega$）。设 $H = G_S$ 为这样的 $S$ 的稳定子。
由于 $[G : H] = \abs{\mathcal O_S}$ 与 $p$ 互素，且 $p^a \mid
\abs G = [G:H]\,\abs H$，得 $p^a \mid \abs H$。反之，取定
$s \in S$，映射 $H \to S$，$h \mapsto hs$ 为单射，且
其像落在 $S$ 中（因 $hS = S$），故 $\abs H \leq \abs S = p^a$。
所以 $\abs H = p^a$。
\end{proof}

\begin{theorem}[西罗 II：支配与共轭]\label{thm:b3:groups:sylow2}
设 $P$ 为一个西罗 $p$-子群，$Q$ 为 $p$-子群中任一个（属于 $G$）。
则 $Q \subseteq gPg^{-1}$，对某 $g \in G$ 成立。特别地，所有
西罗 $p$-子群彼此共轭，且 $P \trianglelefteq G \iff
n_p = 1$。
\end{theorem}

\begin{proof}
令 $Q$ 作用于陪集空间 $X = G/P$ 上，其基数为 $m
\not\equiv 0 \pmod p$。由\cref{thm:b3:groups:pfixed}，用于
$p$-群 $Q$，得 $\abs{X^Q} \equiv m \not\equiv 0 \pmod p$：存在
一个不动陪集 $gP$，即 $QgP = gP$，亦即 $g^{-1}Qg \subseteq
P$。若 $Q$ 本身为西罗子群，则阶的相等使
$Q \subseteq gPg^{-1}$ 化为等式。最后，$P
\trianglelefteq G$ 当且仅当其共轭 $\{gPg^{-1}\}$ ---
由上可知它们恰为\emph{全部}西罗 $p$-子群 --- 归结为
$\{P\}$。
\end{proof}

\begin{theorem}[西罗第三定理：计数]\label{thm:b3:groups:sylow3}
$n_p \equiv 1 \pmod p$，且 $n_p = [G : N_G(P)]$，它整除
$m$。
\end{theorem}

\begin{proof}
记 $\mathrm{Syl}_p$ 为西罗 $p$-子群的集合；$G$ 通过共轭
可递地作用于其上（\cref{thm:b3:groups:sylow2}），
$P$ 的稳定子为正规化子 $N_G(P) \supseteq P$：$n_p = [G :
N_G(P)]$，又 $m = [G:P] = [G:N_G(P)]\,[N_G(P):P]$ 可见 $n_p \mid
m$。

现将该作用限制于 $P$ 并计算不动点。若 $Q \in
\mathrm{Syl}_p$ 被 $P$ 固定，则 $P \subseteq N_G(Q)$；两者
$P$ 与 $Q$ 都是西罗 $p$-子群，二者同在群 $N_G(Q)$ 中，故
在其中共轭（对 $N_G(Q)$ 应用\cref{thm:b3:groups:sylow2}）；
但 $Q \trianglelefteq N_G(Q)$，因此 $Q$ 在其中的共轭仅有它自身：
$P = Q$。于是唯一的不动点为 $P$ 本身，
由\cref{thm:b3:groups:pfixed} 得 $n_p = \abs{\mathrm{Syl}_p}
\equiv \abs{\mathrm{Syl}_p^P} = 1 \pmod p$。
\end{proof}

\begin{method}\label{met:b3:groups:sylowmethod}
分析阶为 $n = p^a m$ 的群时，列出
$m$ 的因数中与 $1$ 模 $p$ 同余者——它们即是
$n_p$ 的候选值。若唯一的候选值为 $1$，则该西罗 $p$-子群
正规。若 $n_p > 1$ 被迫较小，则对
$\mathrm{Syl}_p$ 作共轭作用，得到同态 $G \to S_{n_p}$，
其核较小。再计数元素：两两不同的、阶为\emph{素数} $p$ 的
西罗 $p$-子群两两交平凡，因此共含
$n_p(p-1)$ 个阶恰为 $p$ 的元素；对不同素数的计数
相互重叠，常可导出矛盾
（见\cref{exo:b3:groups:7}）。
\end{method}

\begin{example}\label{ex:b3:groups:pq}
设 $\abs G = pq$，其中 $p < q$ 为素数，且 $p \nmid q - 1$。则
由 $n_q \mid p$ 与 $n_q \equiv 1 \bmod q$ 得 $n_q = 1$（因 $p <
q$）；由 $n_p \mid q$ 与 $n_p \equiv 1 \bmod p$ 得 $n_p = 1$（因 $q
\not\equiv 1 \bmod p$）。设 $P, Q$ 为两个正规西罗子群：$P
\cap Q = \{e\}$（阶互素），故 $\abs{PQ} = pq$
（见\cref{exo:b3:groups:4}），且 $G \cong P \times Q \cong \Z/p\Z
\times \Z/q\Z \cong \Z/pq\Z$，由下文
\cref{prop:b3:groups:direct}，任一阶为 $15$
$33$、$35$、\dots\ 的群都是循环群。例外情形 $p \mid q - 1$
恰多产生一个非交换群——
见周末问题（\cref{pb:b3:groups:1}）。
\end{example}

\begin{example}[西罗子群的完整计数：$S_4$]
\label{ex:b3:groups:s4sylow}
对 $G = S_4$，$\abs G = 24 = 2^3\cdot3$，我们运用这一方法。
\emph{西罗 $3$-子群：}$n_3 \mid 8$，$n_3 \equiv 1 \bmod 3$，故 $n_3
\in \{1, 4\}$；由于 $\langle(123)\rangle$ 与
$\langle(124)\rangle$ 不同，$n_3 = 4$——四个
子群 $\langle(abc)\rangle$，每个对应一个 $3$-元素子集
$\{a, b, c\}$，共包含 $8$ 个三轮换。由西罗第二定理，
它们互相共轭，且共轭同态 $S_4 \to
S_{\mathrm{Syl}_3} \cong S_4$ 在此为同构（其核
包含于 $N = N_G(\langle(123)\rangle)$ 中，该群阶为 $24/4 =
6$，而 $S_4$ 中落在形如 $S_3$ 的 $N$ 内的正规子群必为
平凡群：它只能由偶置换组成，并固定四个
西罗子群，只有 $e$ 满足）。\emph{西罗 $2$-子群：}$n_2 \mid 3$，$n_2
\equiv 1 \bmod 2$：$n_2 \in \{1, 3\}$。子群 $D =
\langle(1234), (13)\rangle$ 的阶为 $8$（二面体群 $D_4$：
即顶点为 $1, 2, 3, 4$ 的正方形的对称群），
不正规（$(12)(1234)(12) = (2134)$ 生成另一个
$4$-轮换子群），故 $n_2 = 3$：三个 $D_4$ 的副本
对应于把 $4$ 个点配成一个
“正方形”的三种方式。注意此次计数的启示：$\abs{S_4} = 24$
给两个西罗子群中任一个不正规留下余地，而两者确实都不正规——
对比阶为 $12$ 的情形，那里计数迫使其中一个正规
（见\cref{pb:b3:groups:1}的第四部分）。
\end{example}

\section{乘积、直积与半直积}

\begin{proposition}[直积的识别]\label{prop:b3:groups:direct}
设 $H, K \trianglelefteq G$，且 $H \cap K = \{e\}$，$HK = G$。
则 $(h,k) \mapsto hk$ 是 $H \times K \to
G$ 的同构。\index{直积}
\end{proposition}

\begin{proof}
对 $h \in H$，$k \in K$，换位子 $hkh^{-1}k^{-1}$ 属于
$K$（将其看作 $(hkh^{-1})k^{-1}$，利用 $K$ 的正规性）且属于
$H$（将其看作 $h(kh^{-1}k^{-1})$）：它等于 $e$，故 $H$ 与 $K$
的元素两两交换，该映射为同态。它是满射
（因 $HK = G$），并且它是单射，因为由 $hk = e$ 得 $h = k^{-1} \in
H \cap K = \{e\}$。
\end{proof}

\emph{两个}因子都正规这一条件最容易不成立：在 $S_3
= \langle (1\,2\,3)\rangle \,\langle(1\,2)\rangle$ 中，两个因子
交平凡且生成，但 $S_3 \not\cong \Z/3\Z \times
\Z/2\Z$。当只有一个因子正规时，正确的概念是：

\begin{definition}\label{def:b3:groups:semidirect}
设 $H$、$K$ 为群，$\varphi \colon K \to
\operatorname{Aut}(H)$ 为同态。\emph{半直
积}\index{半直积} $H \rtimes_\varphi K$ 是
集合 $H \times K$，并赋以如下乘法：
\[
(h, k)\,(h', k') = \bigl(h\,\varphi(k)(h'),\; kk'\bigr).
\]
\end{definition}

\begin{proposition}\label{prop:b3:groups:semidirect}
$H \rtimes_\varphi K$ 是群；$H \times \{e\}$ 是正规
子群，且与 $H$ 同构，$\{e\} \times K$ 是与
$K$ 同构的子群；二者交平凡且生成。反之，若 $G =
NK$，其中 $N \trianglelefteq G$，$K \leq G$，且 $N \cap K = \{e\}$，
则 $G \cong N \rtimes_\varphi K$，其中 $\varphi(k) = (n \mapsto
knk^{-1})$。
\end{proposition}

\begin{proof}
直接验证：结合律归结为 $\varphi(kk') =
\varphi(k)\circ\varphi(k')$，且每个 $\varphi(k)$ 都是同态；
单位元为 $(e,e)$，且 $(h,k)^{-1} =
\bigl(\varphi(k^{-1})(h^{-1}), k^{-1}\bigr)$。投影 $(h,k)
\mapsto k$ 是到 $K$ 上的满同态，核为 $H \times \{e\}$，
因此它是正规的。对于逆命题，每个 $g \in G$ 可
\emph{唯一}地写成 $nk$，其中 $n \in N$，$k \in K$（存在性：$G =
NK$；唯一性：由 $nk = n'k'$ 得 $n'^{-1}n = k'k^{-1} \in N \cap
K$），且
\[
(nk)(n'k') = n\,(kn'k^{-1})\;kk'
\]
上式表明 $nk \mapsto (n, k)$ 将 $G$ 的运算搬到
$N \rtimes_\varphi K$ 的运算上。
\end{proof}

\begin{example}\label{ex:b3:groups:semidirectexamples}
(a) \emph{二面体群}\index{二面体群} $D_n$（$n \geq
3$）是 $2n$ 个对称构成的群，即正 $n$ 边形的全体对称：旋转
构成指数为 $2$ 的正规子群，任一反射生成一个
补子群，用反射共轭一个旋转即得其逆：
$D_n \cong \Z/n\Z \rtimes_\varphi \Z/2\Z$，其中 $\varphi(1) = (x
\mapsto -x)$。
(b) 直线的仿射群 $\{x \mapsto ax + b : a \in
K^\times,\, b \in K\} \cong K \rtimes K^\times$：平移
为正规子群，位似为补子群。
(c) $S_n \cong A_n \rtimes \Z/2\Z$（补子群可取任一
对换）。
(d) 四元数群 $Q_8$ \emph{不是}
两个真子群的半直积：每个非平凡子群都含 $-1$
（见\cref{pb:b3:groups:1}），故任意两个真子群的交
都非平凡。
\end{example}

\section{可解群；\texorpdfstring{$A_n$}{An} 的单性}

\begin{definition}\label{def:b3:groups:derived}
$x, y \in G$ 的\emph{换位子}\index{换位子}为 $[x,y]
= xyx^{-1}y^{-1}$；\emph{导出子群}\index{导出子群}
$D(G)$ 是由全体换位子生成的子群。
\emph{导出列}为 $D^0(G) = G$，$D^{i+1}(G) = D(D^i(G))$，
而 $G$ 称为\emph{可解群}\index{可解群}，若 $D^n(G) =
\{e\}$ 对某个 $n$ 成立。
\end{definition}

\begin{proposition}\label{prop:b3:groups:derived}
$D(G)$ 是正规的（事实上在每个自同构下稳定），$G/D(G)$
为交换群，且对 $N \trianglelefteq G$：$G/N$ 为交换群 $\iff
D(G) \subseteq N$。此外，$G$ 可解当且仅当存在链
$G = G_0 \trianglerighteq G_1 \trianglerighteq \dots
\trianglerighteq G_n = \{e\}$，并使每个 $G_{i+1} \trianglelefteq
G_i$，且每个商 $G_i/G_{i+1}$ 为交换群。可解群的子群与
商群仍可解；反之，若 $N$ 与
$G/N$ 均可解，则 $G$ 亦可解。
\end{proposition}

\begin{proof}
自同构 $\alpha$ 将 $[x,y]$ 映为 $[\alpha x, \alpha y]$：它
置换全体换位子，故保持其生成的子群不变；
共轭是自同构，由此即得正规性。在 $G/D(G)$
中，$\bar x\bar y \bar x^{-1}\bar y^{-1} = \overline{[x,y]} = \bar e$：
商群为交换群。若 $G/N$ 为交换群，则每个 $[x,y] \in
N$，故 $D(G) \subseteq N$；反之，若 $D(G) \subseteq N$，则
$G/N$ 是交换群 $G/D(G)$ 的商，由第三同构
定理，它为交换群。

若 $G$ 可解，则导出列即是这样的链。反之，
给定一条链，由归纳法 $D^i(G) \subseteq G_i$，$G_i/G_{i+1}$
为交换群，故 $D(G_i) \subseteq G_{i+1}$，于是 $D^{i+1}(G) = D(D^i G)
\subseteq D(G_i) \subseteq G_{i+1}$；因此 $D^n(G) = \{e\}$。

遗传性：$D^i(H) \subseteq D^i(G)$（$H \leq G$，由归纳法），
而 $D^i(G/N) = \pi(D^i(G))$，因 $\pi$ 将换位子映到
换位子上；由此即得子群与商群的结论。
扩张性：若 $D^m(G/N) = \{e\}$，则 $D^m(G) \subseteq N$，且
由 $D^n(N) = \{e\}$ 得 $D^{m+n}(G) = D^n(D^m(G)) \subseteq D^n(N)
= \{e\}$。
\end{proof}

\begin{example}\label{ex:b3:groups:solvableexamples}
交换群都可解。$p$-群亦可解，
对阶数归纳：$Z(G) \neq \{e\}$，且 $G/Z(G)$ 是更小的
$p$-群。$S_3$ 与 $S_4$ 可解：$S_4 \trianglerighteq A_4
\trianglerighteq V \trianglerighteq \{e\}$，其中 $V = \{e,
(1\,2)(3\,4), (1\,3)(2\,4), (1\,4)(2\,3)\}$ 为克莱因四元群（
由双对换构成，在 $S_4$ 中正规：它是若干共轭类的并
），商群依次为 $\Z/2\Z$、$\Z/3\Z$、$V$，均为交换群。
在\cref{ch:b3:galois}中，“$n$ 次一般方程
根式可解”在字面上\emph{意味着}“$S_n$ 是
可解群”。由此可见下一定义的重要性。
\end{example}

\begin{definition}\label{def:b3:groups:simple}
称群 $G \neq \{e\}$ 为\emph{单群}\index{单群}，若其
正规子群只有 $\{e\}$ 与 $G$。非交换单群
不可解：$D(G) \trianglelefteq G$ 不是 $\{e\}$
（否则 $G$ 为交换群），于是 $D(G) = G$，导出列
恒定。交换单群恰为 $\Z/p\Z$，$p$
为质数（交换群是单群，当且仅当它没有非平凡真子群，
即由拉格朗日定理，它为质数阶循环群）。
\end{definition}

\begin{lemma}\label{lem:b3:groups:threecycles}
当 $n \geq 3$ 时，$A_n$ 由 $3$-轮换生成；当 $n \geq 5$ 时，
所有 $3$-轮换都\emph{在 $A_n$ 中}共轭。
\end{lemma}

\begin{proof}
$A_n$ 中的元素是偶数个
对换之积；将其两两配对，并使用（复合自右向左）
\[
(a\,b)(c\,d) = (a\,c\,b)(a\,c\,d), \qquad
(a\,b)(b\,c) = (a\,b\,c), \qquad
(a\,b)(a\,b) = e
\]
依次对应不相交、相交与相同的对换对：每一对对换
的乘积都是若干 $3$-轮换之积。

共轭性：$\sigma(a\,b\,c)\sigma^{-1} = (\sigma a\; \sigma b\;
\sigma c)$，故任意两个 $3$-轮换皆可由某个 $\sigma \in
S_n$ 共轭得到。若 $\sigma$ 为奇置换，则将其替换为 $\sigma' = \sigma (d\,e)$，
其中 $d, e$ 是 $\{a, b, c\}$ 之外的两点 --- 由 $n \geq 5$ 可知它们存在；
则 $\sigma'$ 为偶置换，且 $\sigma'(a\,b\,c)
\sigma'^{-1} = \sigma(a\,b\,c)\sigma^{-1}$，因为 $(d\,e)$ 与
$(a\,b\,c)$ 可交换。
\end{proof}

\begin{theorem}[交错群的单性]\label{thm:b3:groups:ansimple}
$A_n$ 当 $n \geq 5$ 时为单群。\index{交错群}
\end{theorem}

\begin{proof}
设 $N \trianglelefteq A_n$，$N \neq \{e\}$。由
\cref{lem:b3:groups:threecycles}，只需证明 $N$
含有 \emph{一个} $3$-轮换：由正规性
及 $3$-轮换在 $A_n$ 中的共轭性，所有 $3$-轮换均属于 $N$，故 $N = A_n$。

对 $\rho \in S_n$，记 $F(\rho) = \{x : \rho(x) \neq x\}$ 为其
\emph{支集}，$f(\rho) = \abs{F(\rho)}$。取 $\sigma \in N
\setminus \{e\}$，使 $f(\sigma)$ \emph{最小}。注意到
非平凡偶置换满足 $f \geq 3$，且 $f(\sigma) =
4$ 对于 $\sigma \in A_n$ 不可能成立，除非 $\sigma$ 为双对换
（$4$-轮换为奇置换）。下证 $\sigma$ 为一个
$3$-轮换。

\emph{情形 A：$\sigma$ 为不相交对换之积}，
不妨设 $\sigma = (a\,b)(c\,d)\cdots$，$f(\sigma) \geq 4$。取 $e'
\notin \{a, b, c, d\}$（由 $n \geq 5$ 可取），令 $\tau =
(c\,d\,e')$，并记
\[
\sigma' = \tau\sigma\tau^{-1}\,\sigma^{-1} \in N
\qquad (\tau\sigma\tau^{-1} \in N \text{ 由正规性}).
\]
由于 $\sigma\tau^{-1}\sigma^{-1} = (\sigma c\;\sigma e'\;\sigma
d) = (d\;\sigma e'\;c)$（利用 $\sigma c = d$，$\sigma d = c$），于是
可得 $\sigma' = (c\,d\,e')(d\;\sigma e'\;c)$。

若 $\sigma e' = e'$（特别地，当 $f(\sigma) =
4$，即 $\sigma = (a\,b)(c\,d)$ 时成立）：则 $(d\,e'\,c) =
(c\,d\,e')$，且 $\sigma' = (c\,d\,e')^2 = (c\,e'\,d)$，是一个
$3$-轮换，位于 $N$ 中，且 $f(\sigma') = 3 < 4 \leq f(\sigma)$
--- 与极小性矛盾。

若 $\sigma e' \neq e'$，则 $\sigma e' \notin \{a, b, c, d,
e'\}$（$\sigma$ 互换 $a,b$ 及 $c,d$，且 $e' \notin \{a,b,c,d\}$
且 $\sigma$ 为单射），故 $\sigma$ 移动六个点 $a, b,
c, d, e', \sigma e'$：$f(\sigma) \geq 6$。另一方面
$\sigma'$ 是两个 $3$-轮换之积，其支集含于 $\{c, d,
e', \sigma e'\}$，满足 $f(\sigma') \leq 4$；且 $\sigma' \neq
e$，因为 $\sigma'(d) = \tau\sigma\tau^{-1}(c) = \tau\sigma(e') =
\sigma e' \neq d$（$\tau$ 固定 $\sigma e' \notin \{c,d,e'\}$）。
故 $\sigma' \in N \setminus\{e\}$，且 $f(\sigma') \leq 4 <
f(\sigma)$，与极小性矛盾。

\emph{情形 B：$\sigma$ 的某个轮换长度 $\geq 3$}，
不妨设 $\sigma(a) = b$，$\sigma(b) = c$，其中 $a, b, c$ 两两不同。若
$\sigma$ 恰为该 $3$-轮换，则证毕。否则
$f(\sigma) \geq 5$（$f(\sigma) = 4$ 且含长度 $\geq3$ 轮换的情形
即奇 $4$-轮换，已排除），故可取 $d, e' \in
F(\sigma) \setminus \{a, b, c\}$。令 $\tau = (c\,d\,e')$，并记
$\sigma' = \tau\sigma\tau^{-1}\sigma^{-1} \in N$。与前相同，
$\sigma' = (c\,d\,e')\,(\sigma c\;\sigma e'\;\sigma d)$ 只移动
以下各点：
\[
M = \{c, d, e'\} \cup \{\sigma c, \sigma d, \sigma e'\}
\subseteq F(\sigma)
\]
（被移动点的像仍被移动：$\sigma(x) \ne x$ 蕴含
$\sigma(\sigma x) \neq \sigma x$，$\sigma$ 为单射）。又
$b \notin M$：五点 $a, b, c, d, e'$ 两两不同，故 $b
\notin \{c, d, e'\}$；若 $b \in \{\sigma c, \sigma d, \sigma
e'\}$，则 $a \in \{c, d, e'\}$（作用 $\sigma^{-1}$，利用
$\sigma a = b$），矛盾。故 $\sigma'$ 固定 $b$，而
$\sigma$ 移动 $b$；且
$F(\sigma') \subseteq F(\sigma)$。最后，$\sigma' \neq e$：
$\sigma^{-1}(c) = b$, $\tau^{-1}(b) = b$, $\sigma(b) = c$,
$\tau(c) = d$，故 $\sigma'(c) = d \neq c$。于是 $\sigma' \in N
\setminus \{e\}$，且 $f(\sigma') \leq f(\sigma) - 1$，
与极小性矛盾。

两种情形均不可能，故 $\sigma$ 为 $3$-轮换。
\end{proof}

\begin{corollary}\label{cor:b3:groups:snnotsolvable}
当 $n \geq 5$ 时，$A_n$ 与 $S_n$ 均不可解，且
$S_n$ 的正规子群仅为 $\{e\}$、$A_n$ 与 $S_n$。
\end{corollary}

\begin{proof}
$A_n$ 是非交换单群，故不可解
（见\cref{def:b3:groups:simple}）；含不可解子群的群
不可解（见\cref{prop:b3:groups:derived}）。设 $N
\trianglelefteq S_n$：则 $N \cap A_n \trianglelefteq A_n$ 等于
$\{e\}$ 或 $A_n$。若 $N \cap A_n = A_n$，则 $A_n \subseteq N$，
且由指数知 $N \in \{A_n, S_n\}$。若 $N \cap A_n = \{e\}$，
则 $N$ 上的投影 $S_n \to S_n/A_n \cong
\Z/2\Z$ 为单射，故 $\abs N \leq 2$；若 $N = \{e, \sigma\}$，
由正规性，$\sigma$ 的共轭类等于
$\{\sigma\}$，即 $\sigma \in Z(S_n)$。但 $Z(S_n) = \{e\}$ 对
$n \geq 3$ 成立：若 $\sigma \ne e$ 将 $a$ 映到 $b \neq a$，取 $c
\notin \{a, b\}$；则 $(b\,c)\sigma(b\,c)^{-1}$ 将 $a$ 映到 $c
\ne b$，故它与 $\sigma$ 不同。于是 $N = \{e\}$。
\end{proof}

\begin{theorem}[若尔当--赫尔德]\label{thm:b3:groups:jordanholder}
每个有限群 $G \neq \{e\}$ 都有一个 \emph{合成
列}\index{合成列}
\[
\{e\} = G_0 \trianglelefteq G_1 \trianglelefteq \cdots
\trianglelefteq G_r = G,
\qquad G_{i}/G_{i-1} \text{ 是单群},
\]
并且 \emph{合成因子} $G_i/G_{i-1}$ 构成的多重集（在
同构意义下）不依赖于所选的合成列。有限
群可解当且仅当其所有合成因子都是
素数阶循环群。\index{若尔当--赫尔德定理}
\end{theorem}

\begin{proof}
\emph{存在性}：对 $\abs G$ 作归纳。若 $G$ 为单群，则取
$\{e\} \trianglelefteq G$。否则取一个极大真正规子群
$N$（子群只有有限多个）；$G/N$ 为单群
（由对应定理，
$G/N$ 的真非平凡正规子群在 $G$ 中的原像
是严格介于 $N$ 与 $G$ 之间的正规子群）。将 $N \trianglelefteq G$ 接到
$N$ 的一个合成列之后。

\emph{唯一性}：对 $\abs G$ 作归纳，$G$ 为单群的情形
显然成立。取两个合成列，其倒数第二项分别为
$M \trianglelefteq G$ 与 $N \trianglelefteq G$（于是 $G/M$、$G/N$
均为单群）。若 $M = N$，则由对 $M$ 的归纳假设得证。
否则，$MN$（它是 $G$ 的正规子群，且严格包含 $M$）
等于 $G$（$M$ 是极大正规子群：若正规子群 $M \subsetneq L
\subsetneq G$ 则其像为
单群 $G/M$ 的真非平凡正规子群，矛盾）。第二同构定理给出
\[
G/M = MN/M \cong N/(M \cap N),
\qquad
G/N = MN/N \cong M/(M \cap N).
\]
记 $K = M \cap N$（$\trianglelefteq G$），取定
$K$ 的一个合成列。于是 $M$ 有两个合成列：
一个是其原有的合成列，另一个是 $K$ 的合成列后接 $K \trianglelefteq
M$（商 $M/K \cong G/N$ 为单群）。由归纳假设（应用于
$M$），$M$ 原有合成列的因子为
$\{\text{因子 } K\} \cup \{G/N\}$；对 $N$ 同理。于是
$G$ 的两个合成列的因子为
\[
\{\text{因子 } K\} \;\cup\; \{\,G/N,\; G/M\,\},
\]
二者为同一个多重集。

可解性：若所有因子都是 $\Z/p_i\Z$，则该列是一条
商均为交换群的链，故 $G$ 可解
（\cref{prop:b3:groups:derived}）。反之，可解群的合成因子
必是可解的（它是某子群的商），且
为单群；可解单群必为交换群（由 $D(G) \ne G$ 可得
$D(G) = \{e\}$），从而为某个 $\Z/p\Z$。
\end{proof}

\begin{remark}\label{rem:b3:groups:jhmeaning}
若尔当–赫尔德定理说，每个有限群都由单群构成，且其部件表
唯一确定 --- 一种群的算术，
其中单群相当于质数，而\emph{各部件如何
粘合}（扩张数据，如半直积中的情形）
取代了单纯的乘法。有限
单群的分类 --- 即循环群 $\Z/p\Z$、交错群 $A_{n \geq
5}$、十六族李型单群，以及 $26$ 个散在单群 --- 是
二十世纪数学的丰碑之一；其证明
分散于约一万页的期刊篇幅中，远远超出
本课程的范围。
\end{remark}

\begin{omfigure}
\begin{tikzpicture}[xscale=1.55, yscale=1.4]
  \tikzset{sg/.style={draw, rounded corners=2pt, inner sep=3pt,
      fill=omDef!6, font=\small}}
  \node[sg, fill=omThm!10] (G)   at (0, 3)    {$D_4$};
  \node[sg] (K1)  at (-1.6, 2)   {$\langle r^2, s\rangle$};
  \node[sg] (R)   at (0, 2)      {$\langle r\rangle$};
  \node[sg] (K2)  at (1.6, 2)    {$\langle r^2, rs\rangle$};
  \node[sg] (S1)  at (-2.4, 1)   {$\langle s\rangle$};
  \node[sg] (S2)  at (-1.2, 1)   {$\langle r^2 s\rangle$};
  \node[sg, fill=omProp!12] (Z) at (0, 1) {$\langle r^2\rangle$};
  \node[sg] (S3)  at (1.2, 1)    {$\langle rs\rangle$};
  \node[sg] (S4)  at (2.4, 1)    {$\langle r^3 s\rangle$};
  \node[sg] (E)   at (0, 0)      {$\{e\}$};
  \draw (G) -- (K1); \draw (G) -- (R); \draw (G) -- (K2);
  \draw (K1) -- (S1); \draw (K1) -- (S2); \draw (K1) -- (Z);
  \draw (R) -- (Z);
  \draw (K2) -- (Z); \draw (K2) -- (S3); \draw (K2) -- (S4);
  \draw (S1) -- (E); \draw (S2) -- (E); \draw (Z) -- (E);
  \draw (S3) -- (E); \draw (S4) -- (E);
\end{tikzpicture}

{\small 二面体群 $D_4 = \langle r,
s \mid r^4 = s^2 = e,\ srs^{-1} = r^{-1}\rangle$ 的十个子群。其中三个
指数为 $2$ 的子群（中间一行）都是正规的，中心也是
$\langle r^2\rangle$（高亮）；四个反射子群
分为两个共轭类，每类两个。自下而上的链
给出合成列，例如 $\{e\} \trianglelefteq \langle
r^2\rangle \trianglelefteq \langle r\rangle \trianglelefteq D_4$：
因子为 $\Z/2\Z, \Z/2\Z, \Z/2\Z$ --- 总是同一个多重集，
这正是若尔当–赫尔德定理所要求的。}
\end{omfigure}

\section{练习}

\begin{exercise}[$\star$]\label{exo:b3:groups:1}
(a) 证明指数为 $2$ 的子群都是正规子群。
(b) 证明若 $[G : H] = 2$，则 $x^2 \in H$ 对每个 $x \in
G$ 成立。
(c) 由此推出 $A_4$ 没有 $6$ 阶子群：拉格朗日定理的
逆命题不成立。\emph{（数一数 $3$-轮换的平方。）}
\end{exercise}

\begin{exercise}[$\star$]\label{exo:b3:groups:2}
证明若 $G/Z(G)$ 是循环群，则 $G$ 是交换群。由此再推出
每个 $p^2$ 阶群都是交换群，并对每个
质数 $p$，给出一个 $p^3$ 阶的非交换群。\emph{（想一想：
$\mathbb F_p$ 上对角元全为单位元的上三角矩阵。）}
\end{exercise}

\begin{exercise}[$\star$]\label{exo:b3:groups:3}
(a) 证明 $\operatorname{Aut}(\Z/n\Z) \cong (\Z/n\Z)^\times$。
(b) 证明内自同构 $\iota_g \colon x \mapsto
gxg^{-1}$ 构成正规子群 $\operatorname{Inn}(G)
\trianglelefteq \operatorname{Aut}(G)$，且
$\operatorname{Inn}(G) \cong G/Z(G)$。
\end{exercise}

\begin{exercise}[$\star\star$]\label{exo:b3:groups:4}
设 $H, K$ 为有限群 $G$ 的子群。
(a) 证明\emph{乘积公式} $\abs{HK}\,\abs{H\cap K} =
\abs H\, \abs K$，通过计数映射 $H \times K
\to HK$，$(h,k) \mapsto hk$ 的纤维。
(b) 证明 $HK$ 为子群当且仅当 $HK = KH$（
两者之一为正规子群时自动成立）。
(c) 若 $H, K \trianglelefteq G$ 且 $H \cap K = \{e\}$，证明
$hk = kh$ 对所有 $h \in H$、$k \in K$ 成立。
\end{exercise}

\begin{exercise}[$\star\star$]\label{exo:b3:groups:5}
（伯恩赛德计数引理）有限群 $G$ 作用于有限
集 $X$。证明轨道的个数为\emph{不动点
个数的平均值}：
\[
\#\{\text{轨道}\} = \frac{1}{\abs G}\sum_{g \in G}
\abs{\operatorname{Fix}(g)},
\qquad \operatorname{Fix}(g) = \{x \in X : g \cdot x = x\},
\]
通过两种方式计数集合 $\{(g,x) : g\cdot x = x\}$。
应用：$\Z/p\Z$（$p$ 为质数）通过旋转作用于
由 $p$ 颗珠子、$a$ 种可用颜色组成的项链上；由此推出费马小定理
$a^p \equiv a \pmod p$。
\end{exercise}

\begin{exercise}[$\star$]\label{exo:b3:groups:6}
利用西罗定理，证明阶为 $15$ 的每个群都是
循环群，且阶为 $45$ 的每个群都是阿贝尔群。
\end{exercise}

\begin{exercise}[$\star\star$]\label{exo:b3:groups:7}
证明阶为 $30$ 的群与阶为 $56$ 的群都不是
单群。\emph{（对于 $30$：若 $n_3 \neq 1$ 且 $n_5 \neq 1$，计数
阶为 $3$ 与 $5$ 的元素。对于 $56$：计数
阶为 $7$ 的元素。）}
\end{exercise}

\begin{exercise}[$\star\star$]\label{exo:b3:groups:8}
(a) 设 $H \leq G$ 的指数为 $n$。证明 $G$ 在
$G/H$ 上的作用给出一个同态 $G \to S_n$，其核 $\bigcap_{g \in
G} gHg^{-1}$ 是 $G$ 中包含于
$H$ 的最大正规子群。
(b) 由此推出：若 $G$ 为有限群，$p$ 是 $\abs G$ 的\emph{最小}质
因数，则指数为 $p$ 的每个子群都是正规子群。
(c) 证明 $G$ 在集合 $X$ 上的每个传递作用
都与陪集空间上的作用同构：存在双射 $X
\to G/G_x$，且与作用相容。
\end{exercise}

\begin{exercise}[$\star\star$]\label{exo:b3:groups:9}
(a) 证明 $D(G)$ 是 $G$ 中使商群为阿贝尔群的最小正规子群，
并证明从 $G$ 到阿贝尔群的每个同态都
唯一地通过\emph{阿贝尔化}
$G^{\mathrm{ab}} = G/D(G)$ 分解。
(b) 计算 $D(S_n)$ 与 $S_n^{\mathrm{ab}}$（$n \geq 2$），以及
$D(Q_8)$ 与 $Q_8^{\mathrm{ab}}$。
\end{exercise}

\begin{exercise}[$\star\star$]\label{exo:b3:groups:10}
设 $G$ 为 $p$-群，$H \subsetneq G$ 为真子群。
证明 $H \subsetneq N_G(H)$（“正规化子会变大”），并由此推出
每个 $p$-群的极大子群都是指数为 $p$ 的正规子群。
\emph{（对 $\abs G$ 归纳，利用 $Z(G) \neq \{e\}$：分别
讨论 $Z(G) \subseteq H$ 与 $Z(G) \not\subseteq H$ 两种情形。）}
\end{exercise}

\begin{exercise}[$\star\star\star$]\label{exo:b3:groups:11}
（$A_5$ 的单性，动手练习）
(a) 证明 $A_5$ 的共轭类的元素个数分别为
$1$、$15$、$20$、$12$、$12$。注意 $S_5$ 中
由 $5$-轮换构成的共轭类的分裂：对一个 $5$-轮换 $\sigma$，比较
$\sigma$ 在 $S_5$ 与 $A_5$ 中的中心化子。
(b) 由此推出 $A_5$ 是单群：正规子群是若干
共轭类的并，包含 $e$，且其阶整除
$60$。
(c) 证明阶为 $60$ 的单群必有 $n_5 =
6$。
\end{exercise}

\begin{exercise}[$\star\star$]\label{exo:b3:groups:12}
（西罗子群的正规化子是自正规化的）设 $P$ 为西罗 $p$-子群，
$G$ 为有限群，且 $H =
N_G(P)$。
(a) 证明 $P$ 是\emph{唯一}的西罗 $p$-子群，
即在 $H$ 中只有一个西罗子群。
(b) 由此推出 $N_G(H) = H$。\emph{（对 $g \in N_G(H)$：$gPg^{-1}$
是西罗 $p$-子群，且位于 $H$ 中，故 $gPg^{-1} = P$。）}
(c) 由此得出：西罗子群的正规化子都不包含于
$G$ 的真正规子群中，并且
包含 $N_G(P)$ 的极大子群是自正规化的。
\end{exercise}

\section{问题：阶不超过 15 的群}

\begin{problem}[{周末问题 --- 小阶
群的分类}]\label{pb:b3:groups:1}
目标是给出同构意义下阶
$\leq 15$ 的群的完全分类，并附完整证明。阶为 $1, 2, 3, 5,
7, 11, 13$ 的群由拉格朗日定理确定（均为循环群），阶为 $4$ 与
$9$ 的群由\cref{thm:b3:groups:pfixed}及下文对 $p^2$ 的分析给出；
其余阶为 $6, 8, 10, 12, 14, 15$。

\medskip
\noindent\textbf{第一部分 --- 工具。}
\begin{enumerate}
  \item 证明每个元素都满足 $x^2 =
        e$ 的群是阿贝尔群；由此推出这样的有限群的阶为
        $2^k$，且同构于 $(\Z/2\Z)^k$。\emph{（将其视为
        $\mathbb F_2$ 上的向量空间。）}
  \item 证明阶为 $p^2$ 的群同构于
        $\Z/p^2\Z$ 或 $(\Z/p\Z)^2$。列出阶为
        $8$ 的阿贝尔群在同构意义下的全部：$\Z/8\Z$、$\Z/4\Z \times
        \Z/2\Z$、$(\Z/2\Z)^3$ --- 证明该列表完备且
        无重复，并\emph{不}借助
        \cref{ch:b3:modules}中的结构定理（按元素的
        最大阶进行讨论）。
  \item 设 $\varphi, \varphi' \colon K \to \operatorname{Aut}(H)$
        为两个作用。证明若 $\varphi' = \varphi \circ
        \alpha$，其中 $\alpha \in \operatorname{Aut}(K)$，则 $H
        \rtimes_{\varphi} K \cong H \rtimes_{\varphi'} K$。
  \item 显式确定 $\operatorname{Aut}(\Z/n\Z)$，其中 $n = 3, 4, 5,
        7$，并证明 $\operatorname{Aut}\bigl((\Z/2\Z)^2
        \bigr) \cong S_3$。
\end{enumerate}

\medskip
\noindent\textbf{第 II 部分 --- 阶为 $2p$ ($6$，$10$，$14$) 与
$pq$ 的群。}
\begin{enumerate}[resume]
  \item 设 $\abs G = 2p$，其中 $p$ 为奇素数。证明 $G$ 有
        一个正规子群 $N = \langle r \rangle$，阶为 $p$，
        以及元素 $s$，其阶为 $2$，且位于 $N$ 之外。
  \item 由此推出 $G \cong \Z/p\Z \rtimes_\varphi \Z/2\Z$，其中
        $\varphi(1) \in \operatorname{Aut}(\Z/p\Z)$ 为对合自同构，
        并得出：$G \cong \Z/2p\Z$ 或 $G \cong
        D_p$；验证这两个群互不同构。
        至此阶为 $6$、$10$、$14$ 的情形得以解决。
  \item 更一般地，设 $\abs G = pq$，其中 $p < q$ 为素数。
        证明若 $p \nmid q-1$，则 $G$ 为循环群
        （见\cref{ex:b3:groups:pq}），并证明若 $p \mid q - 1$，则
        除 $\Z/pq\Z$ 外，\emph{恰有一个}非交换群
        $\Z/q\Z \rtimes \Z/p\Z$（在同构意义下） --- 利用
        第 3 题，以及 $\operatorname{Aut}(\Z/q\Z) \cong
        (\Z/q\Z)^\times$ 为 $q - 1$ 阶循环群，在本阶段不加证明地承认，
        其证明在\cref{ch:b3:galois}中给出（
        $\mathbb F_q^\times$ 的循环性）。由此得出阶为 $15$ 的情形。
\end{enumerate}

\medskip
\noindent\textbf{第 III 部分 --- $8$ 阶群。}设 $G$ 为
        $8$ 阶的非交换群。
\begin{enumerate}[resume]
  \item 证明 $G$ 有元素 $r$，其阶为 $4$（利用
        第 1 题），且 $N = \langle r\rangle$ 为正规子群。
  \item 设 $s \notin N$。证明 $s^2 \in N$
        （见\cref{exo:b3:groups:1}(b)），并证明 $srs^{-1} = r^{-1}$
        （考察 $r$ 在共轭下的可能像，
        它们须为 4 阶元，并排除 $srs^{-1} = r$），并
        证明 $s^2 \in \{e, r^2\}$\emph{（若 $s^2 = r$
        或 $r^3$ 会怎样？为什么 $s^2$ 必与 $s$ 可交换？）}。
  \item 若 $s^2 = e$，证明 $G \cong D_4$。
  \item 若 $s^2 = r^2$，证明乘法表被完全确定，
        所得群即
        \emph{四元数群}\index{四元数群} $Q_8 =
        \{\pm 1, \pm \mathrm i, \pm \mathrm j, \pm \mathrm k\}$,
        $\mathrm i^2 = \mathrm j^2 = \mathrm k^2 = \mathrm i\,
        \mathrm j\,\mathrm k = -1$（取 $r = \mathrm i$，$s =
        \mathrm j$）。验证 $Q_8$ 存在，例如可在
        $GL_2(\C)$ 中，由下式给出
        \[
        \mathrm i \mapsto \begin{pmatrix} \iu & 0\\ 0 & -\iu
        \end{pmatrix},
        \qquad
        \mathrm j \mapsto \begin{pmatrix} 0 & 1\\ -1 & 0
        \end{pmatrix}.
        \]
  \item 证明 $Q_8$ 的每个非平凡子群都含有
        $-1$；由此推出 $Q_8$ 的每个子群都是正规子群，
        $D_4 \not\cong Q_8$（计算 $2$ 阶元的个数），并
        证明 $Q_8$ 不是两个
        真子群的半直积。
\end{enumerate}

\medskip
\noindent\textbf{第 IV 部分 --- $12$ 阶群。}设 $\abs G = 12$，$P_3
\in \mathrm{Syl}_3(G)$, $P_2 \in \mathrm{Syl}_2(G)$.
\begin{enumerate}[resume]
  \item 证明 $n_3 \in \{1, 4\}$，$n_2 \in \{1, 3\}$，并证明若 $n_3
        = 4$ 时必有 $n_2 = 1$\emph{（计算
        $3$ 阶元的个数）}。
  \item 设 $n_3 = 4$。共轭作用在
        $\mathrm{Syl}_3$ 上给出 $\rho \colon G \to S_4$。证明
        $\ker \rho$ 含于每个 $N_G(P_3)$，因而其阶
        整除 $3$，且为平凡群\emph{（它为何不能是
        $3$ 阶的？）}；并证明其像为 $12$ 阶的
        $S_4$ 子群，必为 $A_4$\emph{（指数为 2 的子群都是
        正规的，且包含所有平方元 ---
        见\cref{exo:b3:groups:1}；数出 $S_4$ 中的平方元）}；并
        得出 $G \cong A_4$。
  \item 设 $n_3 = 1$，于是 $G \cong \Z/3\Z \rtimes_\varphi P_2$，
        其中 $\varphi \colon P_2 \to \operatorname{Aut}(\Z/3\Z)
        \cong \Z/2\Z$。逐一列举各种情形：$\varphi$ 平凡
        时得 $\Z/12\Z$ 与 $\Z/6\Z \times \Z/2\Z$；$P_2 =
        \Z/4\Z$ 且 $\varphi$ 满射时，得到
        \emph{广义四元数群} $\mathrm{Dic}_3 = \Z/3\Z \rtimes
        \Z/4\Z$；$P_2 = (\Z/2\Z)^2$ 且 $\varphi$ 满射时，
        在第 3 题的等价意义下，恰得一个
        群 --- 证明它为 $D_6$，例如通过指出一个
        $6$ 阶元与一个类反射的对合元。
  \item 证明 $\Z/12\Z$、$\Z/6\Z \times \Z/2\Z$、$D_6$、$A_4$、
        $\mathrm{Dic}_3$ 两两不同构\emph{（计算
        $2$ 阶元的个数，或利用 $n_3$）}。至此阶为
        $12$ 的情形得以解决。
\end{enumerate}

\medskip
\noindent\textbf{第 V 部分 --- 综合。}
\begin{enumerate}[resume]
  \item 编制分类表：对每个 $n \leq
        15$，列出同构意义下的全部群，并给出
        相应的个数 $1, 1, 1, 2, 1, 2, 1, 5, 2, 2, 1, 5, 1, 2, 1$。
\end{enumerate}

\medskip
\noindent\textbf{第 VI 部分 --- 拓展：阶为 $p^3$ 的群，
$p$ 为奇素数。}第 III 部分中关于 $8$ 阶群的分析有一个漂亮的
奇素数类比，并带有一个真正新的现象。设 $p$ 为
奇素数，$G$ 为 $p^3$ 阶非交换群。
\begin{enumerate}[resume]
  \item 证明 $\abs{Z(G)} = p$，$G/Z(G) \cong
        (\Z/p\Z)^2$\emph{（对中心的循环商迫使群为交换群
        ：见\cref{exo:b3:groups:2}）}，并证明 $D(G) =
        Z(G)$\emph{（证明 $D(G) \subseteq Z(G)$ 时，利用
        $G/Z(G)$ 为交换群；等号成立时，$G$ 为非交换群，
        且 $D(G) \neq \{e\}$）}。由此推出每个换位子
        $[x, y] = xyx^{-1}y^{-1}$ 都是中心元，且其阶
        整除 $p$。
  \item （关键恒等式）设 $x, y \in G$，$z = [y, x]$，
        为中心元。用对 $k$ 的归纳法证明：
        \[
        (xy)^k = x^k y^k z^{k(k-1)/2} .
        \]
        \emph{（把每个 $y$ 移过每个 $x$；每次交叉都会产生
        一个中心因子 $z$。）}
  \item 由此推出：当 $p$ 为奇素数时，映射 $\theta\colon x \mapsto
        x^p$ 是从 $G$ 到 $Z(G)$ 的群同态\emph{（为什么
        $x^p$ 为中心元？为什么 $z^{p(p-1)/2} = e$ 需要 $p$
        为奇数？）}，并由此得出 $G$ 的方次数为 $p$ 或 $p^2$，
        两种情形取决于 $\theta$ 是否
        平凡。
  \item （方次数为 $p$ 的情形）设每个元素都满足 $x^p = e$。
        取 $x, y$，使其陪集生成 $G/Z(G)$，并令 $z =
        [y, x]$。证明 $z \neq e$，$G$ 中每个元素
        都可唯一地写成 $x^ay^bz^c$（$0 \leq a, b, c < p$），并证明
        乘法完全由
        关系 $x^p = y^p = z^p = e$、$z$ 为中心元、$[y, x] =
        z$。验证\emph{海森堡群}
        \[
        H_p = \left\{
        \begin{pmatrix} 1 & a & c\\ 0 & 1 & b\\ 0 & 0 & 1
        \end{pmatrix} : a, b, c \in \mathbb F_p \right\}
        \subseteq GL_3(\mathbb F_p)
        \]
        实现这些关系，且方次数为 $p$（计算
        $(I + N)^p$，其中 $N$ 为严格上三角矩阵，利用
        $N^3 = 0$ 及 $p \geq 3$），每个方次数为 $p$，
        阶为 $p^3$ 的非交换群都同构于 $H_p$。
  \item （方次数为 $p^2$ 的情形）设某个 $r \in G$ 的阶为 $p^2$，
        记 $N = \langle r\rangle$，它是正规的（指数为 $p$：
        见\cref{exo:b3:groups:10}）。证明存在 $s \notin N$，使得
        $s^p = e$\emph{（任取 $t \notin N$；利用
        第 20 题并作修正：$\theta(t) = t^p \in Z(G) \subseteq N$
        --- 证明 $Z(G) = \langle r^p\rangle$ --- 并选取 $a$
        使 $s = tr^{a}$ 满足 $s^p = e$；$p$ 为奇数在何处用到
        ？）}。证明 $srs^{-1} = r^{1+p}$ 在将 $s$ 换为
        某个幂之后成立，并由此得出：恰有\emph{一个}
        阶为 $p^3$ 且方次数为 $p^2$ 的非交换群，
        即 $\Z/p^2\Z \rtimes_\varphi \Z/p\Z$，其中
        $\varphi(1)\colon r \mapsto r^{1+p}$\emph{（利用
        第 3 题；$\operatorname{Aut}(\Z/p^2\Z)$ 为循环群，阶为
        $p(p-1)$，在本阶段不加证明地承认，因而它有唯一的
        $p$ 阶子群）}。
  \item 由此得出计数：对奇素数 $p$，恰有 $5$
        个阶为 $p^3$ 的群（三个交换群，两个非交换群），
        与 $p = 2$ 时一样 --- 但这两个非交换群
        已不再是 $D_4$ 与 $Q_8$。精确指出
        奇 $p$ 的论证在 $p = 2$ 时于何处失效：在
        第 19 题的恒等式中，$z^{k(k-1)/2}$ 当 $k = p = 2$ 时为 $z^1 \neq
        e$，故平方映射不是同态 --- 事实上 $Q_8$
        恰有唯一一个 $2$ 阶元，而方次数为 $4$ 的
        $D_4$ 有五个。
\end{enumerate}

\medskip
\noindent\textbf{第 VII 部分 --- 补充。}
\begin{enumerate}[resume]
  \item 对奇素数 $p$，分别计算两个非交换群中 $p$ 阶元的个数，
        这两个群的阶均为 $p^3$：证明
        $H_p$ 恰有 $p^3 - 1$ 个，而 $M_p =
        \Z/p^2\Z \rtimes \Z/p\Z$ 恰有 $p^2 - 1$ 个
        \emph{（利用第 20 题中的同态 $\theta$：
        先确定其像，再确定其核的阶）}。
        对 $p = 3$ 作数值验证，分别为 $26$ 与 $8$。
        解释为什么这种平方映射论证无法
        区分 $D_4$ 与 $Q_8$，并指出哪一种计数可以
        将它们区分开。
  \item 称整数 $n \geq 1$ 为\emph{循环的}，若每个
        阶为 $n$ 的群都是循环群。证明若 $p^2 \mid n$（对某个素数 $p$），
        或若 $n$ 有素因数 $p < q$，
        $p \mid q - 1$，则 $n$ 不是循环的\emph{（在每种情形下
        给出一个阶为 $n$ 的非循环群，其中第二种情形利用
        第 II 部分）}。由此推出：若 $n$ 为循环的，则
        $\gcd(n, \varphi(n)) = 1$，其中 $\varphi$ 为
        欧拉函数，并对照第 17 题的表检验：
        在 $n \leq 15$ 中，只含一个群的阶恰为
        $n \in \{1, 2, 3, 5, 7, 11, 13, 15\}$，
        正是满足 $\gcd(n, \varphi(n)) = 1$ 的那些数。
\end{enumerate}
\end{problem}
